\documentclass[12pt]{article}
\usepackage[margin=1.15in]{geometry}
\usepackage[T1]{fontenc}
\usepackage{amsmath}
\usepackage{setspace}
\usepackage{hyperref}
\setstretch{1.12}
\setlength{\parindent}{1.4em}
\setlength{\parskip}{0pt}

\title{Graded Evidence for Prehistoric Music:\\Lithophones, Tuning Methods, and Experimental Paradigms}
\author{Flyxion\\Independent Researcher}
\date{October 2026}

\begin{document}
\maketitle

\begin{abstract}
\noindent Claims about the origin of music rest on very thin material, and the strength of a claim depends on which kind of material supports it. This essay arranges the principal evidence for prehistoric music in four grades: instruments whose identity is broadly accepted, objects whose identity is disputed, structures for which an acoustic reading is one hypothesis among several, and indirect inferences from cave acoustics and ethnographic analogy. It then takes the Bruniquel stalagmite structures, which belong to the third grade, as a limiting case, and examines the hypothesis that they were lithophones. The bending-mode physics of free and fixed bars is set out, with a worked calculation that shows how support conditions change both pitch and partial structure, so that fundamental frequency, partial amplitudes, decay, and perceptual discriminability must be considered together. Four tuning procedures based on perceptual feedback are reviewed, and two further procedures are proposed: suspension and rotation of fragments on cords, and sorting of fragments by their partials against a tunable membrane. A set of experimental paradigms follows, each with a stated prediction and a stated outcome that would count against the particular version of the hypothesis that it tests, since a tuned lithophone, an untuned percussion array, acoustic selection of fragments, and coordinated performance are distinct claims. The essay holds throughout that a demonstrated feasibility is not evidence of occurrence, and that the claim which each grade of evidence can bear should be stated at the scope the evidence supports.
\end{abstract}

\section{The problem of a thin record}

Music leaves little. Sound does not fossilize, and the practices that surround it, such as singing, clapping, drumming on perishable surfaces, and rhythmic movement, leave no durable trace unless they happen to involve objects that survive. Every claim about the antiquity of music is therefore a claim about an inference from a small and unrepresentative set of durable objects, and the inference passes through assumptions about what such objects were for. A perforated bone may be a flute or the remains of a carnivore's meal. A ring of broken stalagmites may be a hearth enclosure, a boundary, a ritual setting, or an instrument. The same object supports different claims depending on how much additional structure one is willing to assume.

The aim of this essay is not to settle any of these questions. It is to arrange the available evidence by the kind of claim it can bear, to examine one speculative hypothesis in that light, and to propose experimental paradigms by which parts of that hypothesis could be tested. The hypothesis in question is the Stone Piano Hypothesis, which reads the stalagmite rings of Bruniquel Cave as lithophones, that is, as percussion instruments made of resonant stone \cite{stonepiano}. The hypothesis is presented by its own author as speculative and testable, and as lacking direct evidence in the form of published percussion wear, morphometric data on the fragments, or acoustic measurements of the chamber. That candor is the reason it serves well as a limiting case: it states what it lacks, and the lack defines the work to be done.

\section{Four grades of evidence}

The first grade consists of objects that are generally accepted as instruments. The clearest cases are the flutes from the Swabian Alb in southwestern Germany. Flutes made from mute swan bone and from mammoth ivory at Geissenkl\"osterle have been dated by a high-resolution radiocarbon reexamination to roughly forty-two to forty-three thousand years ago \cite{higham2012}, and a flute made from a vulture radius with five finger holes was recovered at Hohle Fels and dated to the range of thirty-five to forty thousand years, depending on the dating treatment adopted \cite{conard2009}. These objects are worked, have regularly spaced holes made for a purpose that fits sound production, and occur in contexts that also contain figurative art. A claim that deliberate music existed by about forty thousand years ago in this region is well supported, although the claim is about this region and period and does not extend by itself to other places or to earlier times.

The second grade consists of objects whose identity as instruments is disputed. The principal case is a cave bear femur with spaced holes from Divje Babe in Slovenia. An initial radiocarbon date placed it near forty-three thousand years, and a later electron spin resonance estimate placed the bone at fifty to sixty thousand years and indicated that the original dating of that layer was in error. Its discoverer argued that the holes were made by people, while others have argued that carnivore chewing produced them. What the object can bear is a conditional claim: if it is an instrument, then it is a Neanderthal instrument and considerably earlier than the Swabian flutes. The conditional is unlikely to be discharged by further description of the object alone, although microscopic or other examination of how the holes were made could contribute evidence, and in the meantime it should be carried as a conditional in any argument that builds on it.

The third grade consists of structures or assemblages for which an acoustic reading is one hypothesis among several. Bruniquel belongs here. The cave lies in Occitania in southwestern France, was reopened in 1990, and contains structures about 336 meters from the entrance that were dated by uranium series to about 176,500 years ago, with an uncertainty of about two thousand years at the two sigma level \cite{jaubert2016}. They include a larger ring of about 6.7 by 4.5 meters and a smaller one of about 2.2 by 2.1 meters, built from roughly four hundred broken stalagmite pieces with a combined length of about 112 meters and a weight of about 2.2 tonnes. All six structures show signs of fire, including reddened and blackened pieces. The builders were, on the chronology, early Neanderthals. What the structures establish is deliberate construction in a deep cave, the use of fire, and planning. What they do not establish is purpose. Interpretations offered in the literature include a social or symbolic function and a temporary refuge, and the acoustic reading belongs among these as a further candidate.

The fourth grade consists of indirect inferences. Studies of cave acoustics associated with Reznikoff and Dauvois \cite{reznikoff1988} and with Fazenda and colleagues \cite{fazenda2017} have examined whether the locations of painted motifs in Upper Paleolithic caves correlate with acoustic response. Ethnographic lithophones, rock gongs, and resonant stones in several regions show that percussive use of stone is a recurrent human practice. These bodies of evidence establish plausibility and suggest what to measure. They cannot by themselves establish that any particular site was used for music, because an analogy shows what is possible and not what occurred.

\begin{center}
\small
\begin{tabular}{p{0.16\textwidth}p{0.30\textwidth}p{0.42\textwidth}}
\hline
Grade & Representative material & Claim it can bear \\
\hline
First & Swabian flutes (Geissenkl\"osterle, Hohle Fels) & Deliberate music in this region by about forty thousand years ago \\
Second & Divje Babe bone & A conditional claim about an earlier Neanderthal instrument \\
Third & Bruniquel structures & Deliberate deep-cave construction; acoustic purpose is a hypothesis \\
Fourth & Cave acoustics, ethnographic lithophones & Plausibility and measurement targets, not occurrence \\
\hline
\end{tabular}
\end{center}

\section{What the Bruniquel material shows, and what is claimed}

The Stone Piano Hypothesis states that the builders selected, fractured, and arranged stalagmites for their acoustic properties, that the circular arrangement supported coordinated group sound making in a reverberant chamber, and that the deep setting, the firelight, and the absence of domestic debris are compatible with a dedicated performance space \cite{stonepiano}. If correct, it would establish an attestation of musical practice roughly 140,000 years earlier than the Swabian flutes, which would be a much earlier attestation and not a location of music's origin, since earlier practice could have left no durable trace. The paper states three supporting claims, namely that the calcite is acoustically viable, that the fragmentation and arrangement reflect acoustic design, and that the setting is consistent with performance, and it lists falsification criteria that include random length distributions, failure of fragments to produce pitched tones, very short reverberation times, absence of percussion wear, and no compositional difference between selected and unselected stalagmites.

Each supporting claim sits at a different distance from the evidence. The viability of calcite as a ringing material is a matter of laboratory measurement and can be settled independently of Bruniquel. The claim that fragmentation reflects acoustic design is a claim about the Bruniquel fragments and requires morphometric data that are not yet published in the form the hypothesis needs. The claim that the setting suggests performance is the weakest, because the same setting is compatible with several non-acoustic uses, and a setting that is compatible with a hypothesis supports it only to the degree that the setting would be improbable under the alternatives. For these reasons the hypothesis is best held as a program whose first two claims are testable and whose third is interpretive, and the experimental paradigms below are organized accordingly.

The program also contains versions of differing strength that should be kept apart. The strongest is a tuned lithophone, in which fragments were selected or adjusted to produce a set of stable, perceptually distinct pitches. A weaker version is an untuned percussion array, in which stones were struck as sound sources without any stable perceived pitch, long decay, or deliberate clustering of lengths. A third is acoustic selection, in which fragments were chosen for their sound without being tuned relative to one another. A fourth is coordinated performance, in which a group used the structures together in a reverberant space. Stones can serve as instruments in the weaker versions without the properties that the stronger version requires, so an experiment that weakens one version does not exclude the others, and each paradigm below states which version it addresses.

\section{Bending modes of a stone bar}

A stalagmite fragment is, to a first approximation, a tapered rod, and the lowest modes of a rod are bending modes. For a uniform bar of length $L$, Young's modulus $E$, density $\rho$, cross-sectional area $A$, and second moment of area $I$, the Euler--Bernoulli theory gives modal frequencies
\begin{equation}
f_n = \frac{(\beta_n L)^2}{2\pi L^2}\sqrt{\frac{EI}{\rho A}},
\end{equation}
where $\beta_n L$ is determined by the boundary conditions. For a solid circular section of radius $r$ \cite{fletcher} the factor $\sqrt{EI/\rho A}$ equals $(r/2)\sqrt{E/\rho}$. For a bar free at both ends, the first root is $\beta_1 L \approx 4.730$, and the frequencies of the first four modes stand in the ratios $1 : 2.756 : 5.404 : 8.933$. For a bar clamped at one end and free at the other, the first root is $\beta_1 L \approx 1.875$, and the ratios of the first modes are $1 : 6.267 : 17.55$. Two consequences follow. The overtones of a bar are not harmonic multiples of the fundamental, so a struck stone does not give the harmonic series of a string or a pipe, although a dominant partial can still furnish a clear pitch, and fundamental frequency, partial amplitudes, decay, support conditions, and perceptual discriminability all bear on what a listener hears. And the boundary condition changes both the pitch and the overtone structure substantially, which is why a stalagmite that has been broken from its base and rests on a soft support behaves differently from the same stalagmite still attached to the floor.

The magnitude of the effect is easy to illustrate. Take a solid calcite cylinder with $r = 0.025$ m, $E = 75$ GPa, and $\rho = 2600$ kg\,m$^{-3}$, which are the values used in the Stone Piano paper's own worked example. For $L = 0.40$ m the free-free fundamental is about 1490 Hz and the clamped-free fundamental about 235 Hz. For $L = 0.30$ m the free fundamental is about 2660 Hz, and for $L = 0.20$ m it is about 5980 Hz. These figures differ from the value near 320 Hz that the paper's worked example is reported to give, and the discrepancy should be resolved by recomputation before the example is relied upon. The calculation also carries real limitations. Stalagmites taper, their internal structure is layered, the effective modulus of natural speleothem may be well below that of pure calcite, and the length of a fragment is not independent of its radius. The formula is therefore a guide to scaling and to mode ratios and not a prediction of the pitch of any particular fragment. Because frequency scales with the inverse square of length, a reduction of length by ten percent raises the frequency by about twenty-three percent, which is a large audible step and means that removing a small chip from one end is a coarse but usable tuning action.

\section{Tuning and sorting procedures}

The Stone Piano paper describes four procedures by which fragments could have been tuned or selected without abstract acoustic theory \cite{stonepiano}. The first is progressive fracture, in which small chips are removed from one end until the pitch changes in the intended direction. The second is comparative beat matching, in which a fragment is struck in quick succession with a reference and shortened until audible beating disappears. The third is cavity amplification, in which a fragment is held over a natural hollow or cupped hands that act as a resonator reinforcing frequencies near its own. The fourth is assembly-guided sympathetic tuning, in which candidate fragments are tried against a partly built ring and retained if they produce strong sympathetic vibration. All four rely on perceptual feedback and on the physical response of the stone, and none requires measurement.

Two further procedures are proposed here as candidates for experimental investigation. Neither has been shown to have occurred, and each is described so that it can be tested.

The first is suspension and rotation on cords. A bar that rings freely has nodes for its fundamental mode at about $0.224L$ from each end, and a bar suspended by a cord near those positions is minimally damped by its support. Adjusting the position of the cord until the ringing lasts longest therefore locates the nodes without any instrument, and the same adjustment indicates when a fragment is free enough to behave as a free-free bar. A suspended fragment can also be set rotating about the cord, and a small stone point, such as a worked arrowhead, whirled at the end of a cord moves through air at a speed set by the thrower. A bluff or edged body in a flow sheds vortices at a frequency given approximately by $f \approx S\,v/d$, where $v$ is the relative speed, $d$ is the width of the body across the flow, and $S$ is a Strouhal number of about $0.2$ for many shapes. For a point of width $0.02$ m whirled at $10$ m\,s$^{-1}$ the estimate is about $100$ Hz, which lies in the range of an audible low tone. If a point sustains sufficiently coherent periodic airflow, it may produce a measurable tone whose dependence on speed and geometry can be investigated, and shaping or selecting points by such a tone would be a conceivable form of calibration that is entirely distinct from lithophone tuning. The Strouhal estimate does not demonstrate that stable tones occur, or that they are strong enough to be heard or suitable for comparing two objects by beating, and establishing these facts is the first task of the corresponding experiment. The essay treats these as two separable mechanisms, one concerning support and ringing time and the other concerning aerodynamic tone, and the paradigms in the next section test them separately.

The second is sorting by partials against a tunable membrane. A stretched hide, or any membrane on a frame, has its own set of modes, and a circular membrane's lowest modes stand in the ratios $1 : 1.59 : 2.14 : 2.30 : 2.65 : 2.92$, which are also inharmonic. When a struck stone is held close to or lightly in contact with such a membrane, the membrane responds sympathetically if a membrane mode lies near one of the stone's partials. Tension is adjustable, so the membrane can be tuned to a chosen reference, and a collection of fragments can then be sorted into groups according to which membrane mode each excites most strongly. Fragments that excite the same mode share a partial in that region, and fragments that excite a pair of modes simultaneously share more of their overtone structure. This is a way of sorting by partial content using only materials that were available and a response that can be judged by ear or by a visible effect such as the motion of light grains on the surface. Because both stone bars and membranes have inharmonic spectra, the procedure sorts by coincidence of partials and not by harmonic relation in the strict sense, a distinction that matters for what it can demonstrate. The mechanisms involved also need to be separated. Airborne excitation, direct mechanical contact, and loose objects or grains moving on the membrane are different arrangements: contact can alter the very modes being measured, and particle motion can reflect several effects besides frequency coincidence. Each arrangement should be tested separately before the response of a membrane is taken to identify shared stone partials, and the prediction must distinguish shared partials from shared fundamentals and from similar overall spectra.

\section{Experimental paradigms}

The paradigms below are organized by the claim they address. Each states a prediction and an outcome that would count against the hypothesis. Throughout, a successful demonstration of a procedure shows that the procedure is feasible with plausible materials and skills, and it cannot show that the procedure was used at Bruniquel. The inference from feasibility to occurrence requires the further step of comparison with the archaeological record.

\subsection*{Paradigm A: material characterization}

Speleothem analogs of comparable composition and structure are cut or fractured into bars of controlled length and radius, suspended in free-free conditions, struck with a cobble or antler billet, and recorded. Measured quantities are the frequencies and amplitudes of the first several modes, the decay time, and the dependence of the fundamental on length and radius. This paradigm addresses the tuned-lithophone version. The prediction under that version is that fragments in the size range seen at Bruniquel produce clear, sustained, pitched tones and that the fundamental scales with length approximately as the inverse square. A result in which fragments of this size produce short, dull, unpitched sounds regardless of support would count against the tuned version, and would leave the untuned percussion version open, since a stone can be struck as a sound source without a stable pitch or a long decay.

\subsection*{Paradigm B: pitch salience by naive listeners}

Recordings from Paradigm A are played to listeners with no special training, who judge whether each sound has a definite pitch and match it to a reference tone. The prediction is that listeners agree on a pitch for fragments that ring well, and that agreement falls for fragments that are damped or heavily inharmonic. A failure of listeners to extract any stable pitch from well-ringing fragments would indicate that the inharmonic spectrum prevents a usable pitch percept, which would weaken the case that fragments could have been tuned by ear.

\subsection*{Paradigm C: tuning by ear with feedback}

Participants are given a set of fragments longer than needed and a reference sound, and are asked to reduce one fragment by removing chips until it matches the reference or forms a chosen interval with it, using only hearing. The measured outcome is how closely and how quickly the final frequency converges on the target and how the result compares with a physical prediction from the length reduction. The prediction is that perceptual feedback suffices for convergence within a small fraction of a semitone after a modest number of trials. Failure of convergence, or convergence only with measuring instruments, would count against the claim that tuning could have been done without them.

\subsection*{Paradigm D: membrane sorting}

A membrane is mounted on a frame at several tensions, and its modal frequencies are measured. A collection of fragments from Paradigm A is sorted by the membrane mode each excites, and the groups are then compared with groups formed by independent measurement of the fragments' spectra. The prediction is that membrane sorting reproduces spectral grouping with substantial agreement. If the grouping from the membrane is uncorrelated with the measured spectra, the procedure does not work as a sorting device and should be dropped from the list of candidate practices.

\subsection*{Paradigm E: suspension and whirling}

Fragments are suspended from cords at varying positions, and decay time is measured as a function of the suspension point. The prediction is that decay time is maximal near the computed nodal positions and drops measurably as the cord moves away from them. Separately, stone points of varied shape are whirled at controlled speeds, and the experiment first asks whether stable tones occur at all, and with what strength. Only if they do is the dominant frequency measured against speed and width, with the prediction of a frequency proportional to speed and inversely related to width within the range of the Strouhal estimate. A calibration proposal requires two further steps: establishing whether listeners can distinguish differences between objects despite variation in how they are moved, and establishing whether objects that are acoustically matched share the manufacturing properties of interest. The absence of a stable tone, or a flat dependence on speed, would count against the aerodynamic mechanism, while leaving the suspension result unaffected.

\subsection*{Paradigm F: morphometric null models}

The most direct test of the second supporting claim uses the Bruniquel fragments themselves. If lengths and diameters of the fragments can be measured, their distribution is compared with distributions generated by random fracture of stalagmites of the observed sizes, by selection of fallen pieces at random, and by selection to optimize structural stability. Length alone is an insufficient basis for the comparison, as the bar equation makes clear, because frequency depends on thickness, taper, material, and support as well as on length. The test should therefore use the predicted acoustic response of each fragment computed from these quantities, and ask whether the predicted responses depart from those generated by the null processes. Acoustic selection might produce varied lengths that compensate for differences in other dimensions, and regular lengths might instead reflect the requirements of construction, so that neither clustering nor its absence is decisive by itself. The mean fragment length implied by the reported totals, about 112 meters over roughly four hundred pieces, is near 0.28 m, but a mean says nothing about structure and the whole distribution is needed. Absence of departure from the null models in predicted response would count against acoustic selection, and would not by itself bear on the untuned percussion version.

\subsection*{Paradigm G: chamber acoustics and surface wear}

The chamber is modeled, and where permitted measured, for reverberation time and standing-wave locations. The coordinated-performance and tuned versions predict reverberation long enough to sustain overlapping tones, and the prediction of impact-wear marks near antinodes is conditional on the striking practices that experiments show to be plausible, and not a universal expectation. Absent wear becomes informative only where preservation and detection would make it reasonably observable, since fire, deposition, and later disturbance can remove or mimic marks. A very short reverberation time would count against the performance version, and a complete absence of detectable wear on surfaces that would have been struck and that are well preserved would count against percussive use of those surfaces. Surface examination must separate percussion marks from damage produced by fracture, fire, and deposition, and this separation is itself a problem that an experimental series using known impacts on analog stone can help to solve.

\section{A developmental sequence as separable conjectures}

The experimental program above addresses a narrower question than the full hypothesis, which can be read as proposing a developmental sequence in which several practices were connected. Presenting that sequence as a set of separable conjectures preserves the larger idea while allowing each connection to receive its own evidential assessment, and none of the conjectures below is supported by direct evidence at present.

The first conjecture is that skilled movement with weighted cords, of the kind seen in the swinging and whirling of objects in later traditions, supplied a repertoire of embodied control over rhythm and rotation that could be applied to sound making. Its evidential standing is that of an analogy, and a test would examine whether the movements are plausible with materials available in the period and whether they yield sounds that vary systematically with the properties of the swung object, which is what Paradigm E begins to address. The second conjecture is that comparative beat calibration, the adjustment of one sounding object until its beating against a reference disappears, was available as a method, and a test would ask whether untrained participants discover and use it when given only fragments and a reference, as in Paradigm C. The third is that embodied acoustic knowledge, meaning the skills of selecting, striking, listening, and adjusting, could be transmitted between individuals by demonstration and shared practice, without explicit theory. This conjecture can be examined by transmission experiments in which a skilled demonstrator trains novices and the fidelity of the transmitted skill is measured, and it bears on whether the practice could have persisted as a tradition and not remained a one-off event. The fourth is a transition from sounding formations that were already attached in place, such as stalagmites struck where they stood, to constructing an array from detached and arranged fragments. Evidence for it would have to come from the structures themselves, for example from signs that attached formations at the site were also struck or altered, and from comparison of the arrangement with the response of detached pieces.

These conjectures are logically independent. The sequence could be accepted in part, for example the transition from attached to constructed formations without the cord-based movement, and each acceptance would rest on its own tests. Stating them separately also prevents a success on one from being counted as support for the others, which is the form that overreach in this subject most often takes.

\section{Limits of inference}

Three limits deserve stating. The first is that equifinality applies throughout: fragments of similar length may arise from the constraints of stacking and handling, and a ring of stones may be built for reasons that leave no acoustic signature. A positive result in any paradigm must be set against alternative explanations that would produce the same observation. The second is the distinction between feasibility and occurrence already noted. Every procedure proposed here can be shown to work, or shown not to work, and a demonstration that it works enlarges the set of things that were possible without moving any of them into the set of things that happened. The third is that the sample is small. One site, with structures that have been partly disturbed since their making, supports a claim about that site, and extending the claim to Neanderthal music in general requires further material.

These limits do not make the program idle. The first two supporting claims are testable by the paradigms above, and the outcomes are informative whichever way they fall. A negative result in Paradigm A or F would weaken the tuned-lithophone and acoustic-selection versions in particular, and it would leave the untuned percussion version, which does not require stable pitch, long decay, or clustered lengths, to be assessed on other grounds. A positive result would leave the acoustic reading as one of several candidates, strengthened in proportion to how improbable the pattern is under the alternatives. In both cases the claim that results can bear is bounded by the grade of the evidence, which is the arrangement proposed at the outset.

\section{Conclusion}

The history of prehistoric music is a history of inferences from objects, and it is best written by stating for each object what kind of claim it supports. The Swabian flutes support a firm claim about a region and a period. The Divje Babe bone supports a conditional. The Bruniquel structures support a claim about deliberate deep-cave construction and leave purpose open, and the lithophone reading is a hypothesis whose testable parts can be separated from its interpretive part. The experimental paradigms proposed here do not decide the hypothesis. They convert parts of it into measurements with stated outcomes, so that the hypothesis can be refuted, narrowed, or supported by evidence, and so that any later claim made on its behalf can be stated at the scope the evidence warrants.

\begin{thebibliography}{9}
\bibitem{stonepiano} Flyxion, ``The Stone Piano Hypothesis: Neanderthal Lithophones and the Acoustic Interpretation of the Bruniquel Cave Structures,'' 11 March 2026. \url{https://standardgalactic.github.io/calculus/stone_piano.pdf}
\bibitem{jaubert2016} J. Jaubert et al., ``Early Neanderthal constructions deep in Bruniquel Cave in southwestern France,'' \emph{Nature} 534 (2016).
\bibitem{conard2009} N. J. Conard, M. Malina, and S. C. M\"unzel, ``New flutes document the earliest musical tradition in southwestern Germany,'' \emph{Nature} 460 (2009).
\bibitem{higham2012} T. Higham et al., ``Testing models for the beginnings of the Aurignacian and the advent of figurative art and music: the radiocarbon chronology of Gei\ss enkl\"osterle,'' \emph{Journal of Human Evolution} 62 (2012).
\bibitem{reznikoff1988} I. Reznikoff and M. Dauvois, ``La dimension sonore des grottes orn\'ees,'' \emph{Bulletin de la Soci\'et\'e Pr\'ehistorique Fran\c{c}aise} 85 (1988).
\bibitem{fazenda2017} B. Fazenda et al., ``Cave acoustics in prehistory: Exploring the association of Palaeolithic visual motifs and acoustic response,'' \emph{Journal of the Acoustical Society of America} 142 (2017).
\bibitem{fletcher} N. H. Fletcher and T. D. Rossing, \emph{The Physics of Musical Instruments}, 2nd ed., Springer, 1998.
\end{thebibliography}

\end{document}

